% Generated draft: requires mathematical and source review.
\documentclass[11pt]{article}
\usepackage[margin=25mm]{geometry}
\usepackage{fontspec}
\setmainfont{Latin Modern Roman}
\usepackage{amsmath,amssymb,mathtools,unicode-math}
\setmathfont{Latin Modern Math}
\usepackage{enumitem,needspace,xurl,hyperref}
\setlength{\parindent}{0pt}
\setlength{\parskip}{5pt}
\setlength{\emergencystretch}{3em}
\newcommand{\N}{\mathbb{N}}
\newcommand{\Z}{\mathbb{Z}}
\newcommand{\Q}{\mathbb{Q}}
\newcommand{\R}{\mathbb{R}}
\newcommand{\C}{\mathbb{C}}
\newcommand{\F}{\mathbb{F}}
\newcommand{\A}{\mathbb{A}}
\newcommand{\PP}{\mathbb{P}}
\newcommand{\E}{\mathbb{E}}
\newcommand{\eps}{\varepsilon}
\newcommand{\dd}{\,\mathrm{d}}
\begin{document}
\section*{UnsolvedMath ID: 30006160}
\textbf{Source problem code:} catalogue entry
\subsection{Definitions and mathematical statement}
\textbf{Unrecoverable statement.} The supplied record for UnsolvedMath ID \(30006160\), source code \texttt{catalogue entry}, contains the isolated word ``conjecture'' followed by table-of-contents entries. No mathematical assertion, definitions, hypotheses, domains, boundary conditions, or quantifiers are recoverable from this fragment. Consequently, no precise conjecture can be stated from the supplied evidence.

The contents fragment lists: ``conjecture'' (p.~100); Eyal Kaplan, \emph{Consistency results regarding the Ketonen order and the Lipschitz order} (p.~101); Matthew Foreman, \emph{Filter Games and Ineffable Cardinals} (p.~103); Moti Gitik, \emph{Some new ways of blowing up the power of a measurable cardinal} (p.~104); Filippo Calderoni (joint with Luca Motto Ros), \emph{Idealistic equivalence relations remastered} (p.~105); and Elliot Glazer, \emph{A hypertalk} (p.~106). These entries do not identify an intended conjecture or its proposer.

The registry labels the record ``open'' and its statement status ``unrecoverable.'' The open label is supplied metadata, not a verified mathematical status for an identifiable proposition.
\subsection{Short English statement}
This record preserves a table-of-contents fragment from a 2025 set theory workshop report. It does not supply enough information to formulate a mathematical problem. The intended conjecture, its assumptions, and its attribution require source recovery and human review.
\subsection{Sources}
\begin{itemize}
\item Supplied local UnsolvedMath registry record: permanent ID \(30006160\), source code \texttt{catalogue entry}, title \emph{Descriptive-Set-Theory Table-of-Contents Fragment}; collection: \emph{Oberwolfach Reports Open Problems}. No proposer is specified; the supplied proposed year is 2025.
\item Supplied report citation: \emph{Set Theory} (2025), Oberwolfach Reports, pp.~75--140, DOI \texttt{10.4171/owr/2025/2}. \url{https://doi.org/10.4171/owr/2025/2}.
\item Supplied literature source: \emph{Oberwolfach Report 2/2025}. \url{https://content.ems.press/assets/public/full-texts/serials/owr/22/1/14299082/online/10.4171-owr-2025-2.pdf}.
\item Literature assessment dated \texttt{2026-08-22}: the supplied curation labels the record open but reports that the fragment contains no recoverable testable proposition. This is supplied dated curation, not an independently verified literature assessment. Its background explicitly requests independent verification of time-sensitive status and cited claims.
\item The supplied statement-verification field reports that the official report confirms the extraction is table-of-contents debris associated with the word ``conjecture.'' This verification claim is reproduced as supplied provenance; no independent source consultation is asserted here.
\end{itemize}
\end{document}
